\documentclass[conference]{IEEEtran}
\IEEEoverridecommandlockouts
\usepackage[T1]{fontenc}
\usepackage{cite}
\usepackage{url}
\usepackage{amsmath,amssymb,amsfonts}
\usepackage{graphicx}
\usepackage{textcomp}
\usepackage{xcolor}
\usepackage{booktabs}
\begin{document}

% Stichpunkt-Liste mit wenig Abstand
\newenvironment{points}{\begin{itemize}\setlength{\itemsep}{1pt}\setlength{\parskip}{0pt}}{\end{itemize}}

% Stand: Fortschrittsbericht 1, ueberarbeitet nach Feedback (Punkte 1-12)
%  - Kap. I + II ausformuliert und belegt; Kap. III-V stichwortartig
%  - Frameworks nur in III-E "Implementation Details"

\title{Rules or LLM Agents? Real-Time Behaviour\\ Selection for an Emotion-Responsive Robot}

\author{\IEEEauthorblockN{Alexander Scheidl}
\IEEEauthorblockA{UAS Technikum Wien}
\and
\IEEEauthorblockN{Victor Swekis}
\IEEEauthorblockA{UAS Technikum Wien}
}

\maketitle

\begin{abstract}
\textit{Outline (written last):}
\begin{points}
    \item Problem: social robots must react to a person's emotional state within about one second; LLM agents promise more flexible reactions than hand-crafted rules, but it is unknown whether they meet this timing on robot hardware.
    \item Method: controlled comparison of a rule-based controller, LLM-agent variants, a hybrid and a distilled small model, with identical perception and actions, on the onboard computer and on an edge GPU.
    \item Result and conclusion: latency, resources, validity, consistency and appropriateness; answer to the research question.
\end{points}
\end{abstract}

%%%%%%%%%%%%%%%%%%%%%%%%%%%%%%%%%%%%%%%%%%%%%%%%%%%%%%%%%%%%%%%%%%
\section{Introduction}

Social robots are robots designed to interact with people in a socially
meaningful way, communicating through speech, gaze, gesture and
expression \cite{Fong2003Survey}. They are used as companions, assistants
and tutors, and they belong to a growing professional service-robot market
whose global sales rose by 24\,\% to almost 250\,000 units in 2025
\cite{IFR2026Service}.

Whether such a robot succeeds depends less on its mechanics than on the
quality of the interaction, which is the subject of human--robot interaction
(HRI) research \cite{Goodrich2007HRI}. A central requirement is affective
responsiveness: the robot has to recognise how a person feels and react to it
appropriately. Emotion recognition from faces, body and voice has matured
to the point where it runs in real time during interaction with physical
robots \cite{Spezialetti2020Emotion}. The reaction also has to come in time: users
prefer robot responses after about one second, and a robot should react
within two \cite{Shiwa2008Respond}. Even a robot with a strongly reduced
expressive channel can meet this requirement: observers rate the low-cost
Reachy Mini as warmer and more sociable when it expresses positive affect
through head and antenna motion alone \cite{Rogel2026Reachy}.

Between recognition and expression sits the \emph{behaviour-selection layer},
which maps the perceived state of a person to the robot's reaction. In
robotics practice, this layer is predominantly hand-crafted, typically as a
finite-state machine or a behaviour tree
\cite{Ghzouli2023BTSM,Iovino2022BTSurvey}. Learned policies, mainly reinforcement
learning, have been explored, but they require a hand-designed reward and
extensive interaction data and so far remain confined to research studies
\cite{Akalin2021RL}. Hand-crafted rules are fast and predictable, but every
situation has to be anticipated, and repeated situations lead to identical,
mechanical reactions.

Large language models (LLMs) offer a different approach. Used as
tool-calling agents \cite{Yao2023ReAct}, they select an action from a textual
description of the situation and the interaction history, and users prefer
LLM-powered robots for social tasks \cite{Kim2024LLMHRI}. LLM inference,
however, takes hundreds of milliseconds to seconds and needs considerable
compute and memory \cite{Abstreiter2025Edge}, and even strong models do not
always produce valid tool calls \cite{Patil2025BFCL}. On the embedded
computers of small robots, this conflicts directly with the one-second
requirement. Existing LLM-based HRI studies evaluate perceived quality with
models running off the robot \cite{Kim2024LLMHRI}; a
controlled comparison with rule-based behaviour selection under real-time and
on-device constraints is missing.

\textbf{Research question.}
\begin{quote}
\emph{Can an LLM agent running locally on robot or edge hardware select
emotion-responsive behaviour within the one-second response window, and what
does it offer over a rule-based controller?}
\end{quote}

\textbf{Contribution.} This paper provides empirical grounding for the
choice between hand-crafted and LLM-based behaviour selection in affective
HRI. It isolates the behaviour-selection layer as the only variable, keeping
perception, action repertoire and motion execution fixed, and measures what
the timing requirement of social interaction permits on hardware a small
social robot can carry. From this, it derives under which conditions an LLM
agent, a hybrid or a distilled small model is worth its cost compared with
rules.

\textbf{Verification.} The first part of the research question is answered
by the reaction latency of each configuration relative to the one- and
two-second thresholds of \cite{Shiwa2008Respond}, measured on the robot's
onboard computer and on an edge GPU. The second part is answered by the
validity, run-to-run consistency and appropriateness of the selected
reactions against an independent annotation, and by the resource demand of
each configuration (Section~\ref{sec:eval}).

%%%%%%%%%%%%%%%%%%%%%%%%%%%%%%%%%%%%%%%%%%%%%%%%%%%%%%%%%%%%%%%%%%
\section{State of the Art}

Behaviour selection for social robots can be realised with four families of
methods: hand-crafted control structures, learned policies, LLM-based
selection, and combinations of them. This section compares them by the
requirements derived above, namely reaction within about one second,
appropriate and non-mechanical reactions, and operation on the robot's own
hardware. Table~\ref{tab:sota} summarises the comparison.

\subsection{Hand-Crafted Control Structures}
Finite-state machines (FSMs) have long been the standard for modelling robot
behaviour; behaviour trees (BTs) organise the same logic hierarchically and
are increasingly used in open-source robotics
\cite{Ghzouli2023BTSM,Iovino2022BTSurvey}. Both are executed in
microseconds, are deterministic and can be verified before deployment
\cite{Iovino2022BTSurvey}. Their limitation is coverage: every situation and
every variation of a reaction has to be designed by hand, so unforeseen
situations are not handled and repeated situations produce identical
reactions.

\subsection{Learned Policies}
Reinforcement learning (RL) lets a social robot adapt its behaviour to a
user, with rewards derived from task success, user feedback or affective
signals \cite{Akalin2021RL}. Once trained, an RL policy is as fast as a rule
table. However, it needs a reward function that captures what an appropriate
social reaction is, and it learns from many interactions with real users;
the reviewed studies therefore address narrow tasks in controlled settings
\cite{Akalin2021RL}. For open emotion-responsive behaviour without a training
population, RL is not applicable, which is why it is not used here.

\subsection{LLM-Based Behaviour Selection}
LLMs select actions in two ways. As \emph{agents}, they interleave reasoning
with tool calls until an action is chosen \cite{Yao2023ReAct}. In
\emph{single-call} selection, the model
answers once with an output constrained to a schema or grammar, which
guarantees syntactically valid actions at almost no extra cost
\cite{Willard2023Guided}. In HRI, LLM-powered robots are preferred for
social tasks, although verbose or slow responses reduce this preference
\cite{Kim2024LLMHRI}. Cloud models add a network round trip and send
interaction data to an external service. Locally, latency depends strongly on model size and
hardware: on single-board computers, generation takes seconds, while edge
GPUs reduce it considerably \cite{Abstreiter2025Edge}. Small models can be
specialised for one task by knowledge distillation from a larger teacher
\cite{Hinton2015Distill}; models with a few hundred million parameters
trained for function calling now exist.\footnote{FunctionGemma 270M,
\url{https://huggingface.co/google/functiongemma-270m-it}}

\subsection{Hybrid Architectures}
Dual-process architectures combine a fast, automatic layer with a slower,
deliberative one, following the System~1/System~2 distinction of human
cognition \cite{Umbrico2022Miriam}. Applied to behaviour selection,
time-critical reactions can be handled by rules while the LLM handles
situations where context matters, so that the latency of the LLM only
affects reactions that tolerate it.

\begin{table}[t]
\caption{Behaviour-selection methods against the requirements of
emotion-responsive HRI on robot hardware ($+$ met, $\circ$ partly, $-$ not met).}
\label{tab:sota}
\centering
\setlength{\tabcolsep}{3pt}
\begin{tabular}{@{}lccccc@{}}
\toprule
Method & $<$1\,s & Valid & Flexible & No training & On-device \\
\midrule
FSM / BT \cite{Ghzouli2023BTSM}      & $+$ & $+$ & $-$ & $+$ & $+$ \\
RL policy \cite{Akalin2021RL}        & $+$ & $+$ & $\circ$ & $-$ & $+$ \\
LLM agent \cite{Yao2023ReAct}        & $-$ & $\circ$ & $+$ & $+$ & $\circ$ \\
LLM single call \cite{Willard2023Guided} & $\circ$ & $+$ & $+$ & $+$ & $\circ$ \\
Hybrid \cite{Umbrico2022Miriam}      & $+$ & $+$ & $\circ$ & $+$ & $\circ$ \\
Distilled SLM \cite{Hinton2015Distill} & $\circ$ & $\circ$ & $\circ$ & $\circ$ & $+$ \\
\bottomrule
\end{tabular}
\end{table}

\subsection{Research Gap}
Table~\ref{tab:sota} shows that no single method meets all requirements.
Rules meet the timing but not the flexibility, LLMs the flexibility but not
the timing on robot hardware, and RL is ruled out by its training demand.
Hybrid architectures and distilled small models are proposed as a way out,
but they have not been measured against rules and full LLM agents under
identical conditions. What is missing is therefore not a new selection
method, but a controlled comparison of these candidates on the hardware a
small social robot can use. Answering the research question requires
exactly this comparison, which this paper provides.

%%%%%%%%%%%%%%%%%%%%%%%%%%%%%%%%%%%%%%%%%%%%%%%%%%%%%%%%%%%%%%%%%%
\section{Materials and Methods}

The method keeps everything except the behaviour-selection layer fixed, so
that differences in the results can be attributed to that layer alone.

\subsection{Platform and Hardware}
\begin{points}
    \item Reachy Mini \cite{Rogel2026Reachy}: 6-DoF head, body rotation, two antennas, camera.
    \item Phase 1: simulation on a laptop, to develop and calibrate all configurations.
    \item Phase 2a: Reachy Mini Wireless, everything on the onboard Raspberry Pi CM4 (4\,GB RAM, no GPU) -- the hardware a small social robot carries.
    \item Phase 2b: LLM inference offloaded to an NVIDIA Jetson AGX Orin edge GPU next to the robot -- local, without cloud, but with GPU acceleration \cite{Abstreiter2025Edge}.
\end{points}

\subsection{System Architecture}
\begin{points}
    \item Perception: face position, distance, emotion class, hand gestures.
    \item Event filter: only relevant, debounced changes reach the behaviour selection.
    \item Motion layer (50\,Hz): head keeps tracking the face; expressions are overlaid without blocking, so a slow selection never stops tracking.
    \item Shared action repertoire: seven emotional postures, wave, nod, no action.
\end{points}

\subsection{Compared Configurations}
Each configuration represents one method family of Section~II that is
applicable without a training population (Table~\ref{tab:sota}):
\begin{points}
    \item R -- rule table: representative of hand-crafted FSM/BT control, the dominant practice (baseline) \cite{Ghzouli2023BTSM}.
    \item A-TC -- tool-calling agent, one tool per action: the standard LLM-agent paradigm \cite{Yao2023ReAct}.
    \item A-SO -- single LLM call with schema-constrained output: removes the agent loop and guarantees valid actions \cite{Willard2023Guided}.
    \item H -- hybrid: rules for time-critical events (person appears, gesture, distance), A-SO for the rest \cite{Umbrico2022Miriam}.
    \item D -- distilled small model: a 270M-parameter function-calling model (FunctionGemma) fine-tuned on decisions of a larger teacher model \cite{Hinton2015Distill}, small enough for the onboard computer.
    \item Not compared: RL (no training population, Section~II-B) and cloud LLMs (network latency, interaction data leave the robot).
\end{points}

\subsection{Evaluation Design}
\label{sec:eval}
\begin{points}
    \item Input A -- scripted event sequence (36 events): decision only, without perception noise.
    \item Input B -- recorded stimulus video with scripted emotions and gestures, replayed in real time: end to end.
    \item No existing benchmark fits: HRI datasets such as UE-HRI annotate user engagement \cite{BenYoussef2017UEHRI}, not the appropriateness of a robot reaction to an emotional state; stimuli are therefore recorded and annotated by the authors.
    \item Annotation: two authors independently define the set of acceptable reactions per event before seeing results; a large cloud LLM labels the same events from the symbolic event log (no video, so no personal image data leaves the system) as a third annotator \cite{Gilardi2023Annotation}; agreement reported as Cohen's $\kappa$ \cite{McHugh2012Kappa}.
    \item Ten interleaved runs per configuration and input, per hardware phase.
\end{points}

\subsection{Metrics}
\begin{points}
    \item Latency: decision latency (event to action) and end-to-end latency (camera frame to motor command); median, 95th percentile and share of reactions above 1\,s and 2\,s \cite{Shiwa2008Respond}.
    \item Validity: share of outputs that are an executable action with valid arguments, checked structurally as in the Berkeley Function Calling Leaderboard \cite{Patil2025BFCL}.
    \item Consistency: agreement of the selected action across the ten runs per event, Fleiss' $\kappa$ \cite{McHugh2012Kappa}.
    \item Appropriateness: share of reactions within the annotated acceptable set, and Cohen's $\kappa$ between configuration and annotation \cite{McHugh2012Kappa}.
    \item Resources: CPU, RAM and (on the Jetson) GPU load; whether perception still holds 30\,fps next to the selection.
\end{points}

\subsection{Implementation Details}
\begin{points}
    \item Robot control via the Reachy Mini SDK\footnote{\url{https://github.com/pollen-robotics/reachy_mini}}; simulation in MuJoCo\footnote{\url{https://mujoco.org}}.
    \item Perception with MediaPipe face landmarks and gesture recognition.\footnote{\url{https://ai.google.dev/edge/mediapipe}}
    \item A-TC built with LangChain\footnote{\url{https://github.com/langchain-ai/langchain}}; local inference via Ollama\footnote{\url{https://ollama.com}}; teacher and agent models from the Qwen2.5 family.
    \item Agent receives a textual state description plus the last five events, no images; tools enqueue the action directly (no second LLM call).
\end{points}

%%%%%%%%%%%%%%%%%%%%%%%%%%%%%%%%%%%%%%%%%%%%%%%%%%%%%%%%%%%%%%%%%%
\section{Experimental Results}
\textit{Planned content -- measurements pending.}
\begin{points}
    \item Latency per configuration on CM4 and Jetson against the 1\,s/2\,s thresholds.
    \item Resource demand; does the selection fit on the robot next to perception?
    \item Validity, consistency and appropriateness of the reactions.
    \item Discussion: trade-off between latency/cost and flexibility; limitations.
\end{points}

%%%%%%%%%%%%%%%%%%%%%%%%%%%%%%%%%%%%%%%%%%%%%%%%%%%%%%%%%%%%%%%%%%
\section{Summary and Outlook}

This paper examines whether a locally running LLM agent can select
emotion-responsive behaviour for a social robot within the one-second
response window, and what it offers over rules. It compares a rule table, a
tool-calling agent, a single-call LLM, a hybrid and a distilled small model
with identical perception and actions on the Reachy Mini's onboard computer
and on an edge GPU, measuring latency, validity, consistency,
appropriateness and resource demand.

Future work will complete the measurements and test the resulting
recommendation in a user study with established questionnaires,
to check whether the measured differences are also perceived by users.

%%%%%%%%%%%%%%%%%%%%%%%%%%%%%%%%%%%%%%%%%%%%%%%%%%%%%%%%%%%%%%%%%%
{\small
\bibliographystyle{IEEEtranS}
\bibliography{references}
}

\end{document}
